% figures05.tex -- Chapter 5, "Measurement of Segments".
% All figures in this chapter are number-line / segment diagrams: labelled
% points on straight lines, with the unit segment shown below.  Every point
% position below is exact arithmetic on the unit u -- no point is placed by
% eye.  Proportions read off the iPad scans (scan05 pp.18-24, scan06 pp.6-7)
% and the photo PDF (PDF pp.87-97); see tools/constraints/ch05.py.
%
% Fig. 5-1 (a pictorial vignette of a man pacing off a building) is NOT
% redrawn -- see chapters/ch05-UNCERTAIN.md.
%
% Loaded here rather than in ch05.tex so that this file stands alone for
% tools/check_labels.py, whose driver inputs figures05.tex by itself.
\usetikzlibrary{decorations.pathreplacing}

% ---------------------------------------------------------------- Fig. 5-2
% Unit CD laid off from A: A1,A2,A3,A4 at 1..4 units; B between A3 and A4
% (measured at 0.36 of that gap, giving AB about 3.36 units).
\newcommand{\FIGVTWO}{%
\begin{bkfigure}[0.92\linewidth]
\begin{tikzpicture}[x=1.75cm,y=1cm,scale=1.257]
  \draw[given] (-0.20,0) -- (4.40,0);
  \foreach \x/\lb in {0/A,1/{A_1},2/{A_2},3/{A_3},3.36/B,4/{A_4}}{%
    \dt{\x,0}
    \node[vlab,above] at (\x,0) {$\lb$};}
  \node[vlab,right] at (4.40,0) {$l$};
  % unit segment CD, drawn below, spanning 0.52u to 1.52u.  Drop measured
  % from the book plate: gap between the two lines is 0.35 of the unit.
  \draw[given] (0.52,-0.61) -- (1.52,-0.61);
  \dt{0.52,-0.61} \dt{1.52,-0.61}
  \node[vlab,above,xshift=-2pt] at (0.52,-0.61) {$C$};
  \node[vlab,above,xshift=2pt] at (1.52,-0.61) {$D$};
  \node[slab,above] at (1.02,-0.63) {1};
\end{tikzpicture}
\figcap{5--2}
\end{bkfigure}}

% ---------------------------------------------------------------- Fig. 5-3
% The unit is too large: B lies between A and A1.  Re-measured off the book
% plate: B sits at 0.32 of AA1 (not 0.45), and the line runs on to 1.71 u.
\newcommand{\FIGVTHREE}{%
\begin{bkfigure}[0.62\linewidth]
\begin{tikzpicture}[x=1.75cm,y=1cm,scale=1.300]
  \draw[given] (-0.26,0) -- (1.71,0);
  \foreach \x/\lb in {0/A,0.32/B,1/{A_1}}{%
    \dt{\x,0}
    \node[vlab,above] at (\x,0) {$\lb$};}
\end{tikzpicture}
\figcap{5--3}
\end{bkfigure}}

% ----------------------------------------------------------- Figs. 5-4/5-5
% Archimedes' integer n = 4.  5-4: B = A4 exactly.  5-5: B between A3 and A4
% (measured at 0.18 of that gap).  In both plates the drawn line STARTS at A
% -- there is no stub to the left of A, unlike Fig. 5-2.
\newcommand{\FIGVFOURFIVE}{%
\begin{bkfigure}[0.94\linewidth]
\begin{tikzpicture}[x=1.75cm,y=1cm,scale=1.274]
  \draw[given] (0,0) -- (4.40,0);
  \foreach \x/\lb in {0/{A=A_0},1/{A_1},2/{A_2},3/{A_3},4/{A_4=B}}{%
    \dt{\x,0}
    \node[vlab,above] at (\x,0) {$\lb$};}
  \draw[given] (0.52,-0.61) -- (1.52,-0.61);
  \dt{0.52,-0.61} \dt{1.52,-0.61}
  \node[vlab,above,xshift=-2pt] at (0.52,-0.61) {$C$};
  \node[vlab,above,xshift=2pt] at (1.52,-0.61) {$D$};
\end{tikzpicture}
\figcap{5--4}
\end{bkfigure}%
\begin{bkfigure}[0.94\linewidth]
\begin{tikzpicture}[x=1.75cm,y=1cm,scale=1.274]
  \draw[given] (0,0) -- (4.40,0);
  \foreach \x in {0,1,2,3,3.18,4}{\dt{\x,0}}
  \foreach \x/\lb in {0/A,1/{A_1},2/{A_2},4/{A_4}}{%
    \node[vlab,above] at (\x,0) {$\lb$};}
  % A3 and B sit 0.18 u apart, as in the book; the labels are nudged apart
  % by 1.5pt each so they read as "A_3 B" rather than colliding.  The DOTS
  % stay at the exact positions 3 and 3.18 -- only the text moves.
  \node[vlab,above,xshift=-1.5pt] at (3,0) {$A_3$};
  \node[vlab,above,xshift=1.5pt] at (3.18,0) {$B$};
\end{tikzpicture}
\figcap{5--5}
\end{bkfigure}}

% ----------------------------------------------------------- Figs. 5-6/5-7
% 5-6: the unit interval CD divided into tenths (9 interior marks).
% 5-7: AB measured as about 2.3 units; the tenth-scale unit CD alongside.
\newcommand{\FIGVSIXSEVEN}{%
\begin{bkfigure}[0.90\linewidth]
\begin{tikzpicture}[x=7.5cm,y=1cm,scale=1.300]
  \draw[given] (0,0) -- (1,0);
  \dt{0,0} \dt{1,0}
  \node[vlab,above,xshift=-2pt] at (0,0) {$C$};
  \node[vlab,above,xshift=2pt] at (1,0) {$D$};
  \foreach \k in {1,...,9}{\draw[tick] ({\k/10},-0.085) -- ({\k/10},0.085);}
\end{tikzpicture}
\figcap{5--6}
\end{bkfigure}%
\begin{bkfigure}[0.90\linewidth]
\begin{tikzpicture}[x=1.9cm,y=1cm,scale=1.300]
  \draw[given] (0,0) -- (2.30,0);
  \foreach \x/\lb in {0/A,1/{A_1},2/{A_2},2.30/B}{%
    \dt{\x,0}
    \node[vlab,above] at (\x,0) {$\lb$};}
  % the unit CD alongside, itself divided into tenths
  \draw[given] (2.64,0) -- (3.64,0);
  \dt{2.64,0} \dt{3.64,0}
  \node[vlab,above,xshift=-2pt] at (2.64,0) {$C$};
  \node[vlab,above,xshift=2pt] at (3.64,0) {$D$};
  \foreach \k in {1,...,9}{\draw[tick] ({2.64+\k/10},-0.075) -- ({2.64+\k/10},0.075);}
\end{tikzpicture}
\figcap{5--7}
\end{bkfigure}}

% ---------------------------------------------------------------- Fig. 5-8
% Theorem 5-2: B'' chosen on the same side of A as B with AB'' = A'B'.
% B'' drawn between A and B (measured at 0.77 of AB).  AB is drawn the same
% length as A'B' (the theorem's given |AB| = |A'B'|), A sits directly under
% A', and the two rows are 0.17 of the segment length apart -- all measured
% off the book plate.
\newcommand{\FIGVEIGHT}{%
\begin{bkfigure}[0.72\linewidth]
\begin{tikzpicture}[x=4.6cm,y=1cm,scale=1.300]
  \draw[given] (0,0.80) -- (1,0.80);
  \dt{0,0.80} \dt{1,0.80}
  \node[vlab,above] at (0,0.80) {$A'$};
  \node[vlab,above] at (1,0.80) {$B'$};
  \draw[given] (0,0) -- (1,0);
  \foreach \x/\lb in {0/A,1/B}{%
    \dt{\x,0}
    \node[vlab,above] at (\x,0) {$\lb$};}
  \dt{0.77,0}
  \node[vlab,above,yshift=1.2pt] at (0.77,0) {$B''$};
\end{tikzpicture}
\figcap{5--8}
\end{bkfigure}}

% ---------------------------------------------------------------- Fig. 5-9
% Theorem 5-3: E on the same side of C as D with AB = CE; D between C and E
% (measured at 0.66 of CE).
\newcommand{\FIGVNINE}{%
\begin{bkfigure}[0.72\linewidth]
\begin{tikzpicture}[x=4.6cm,y=1cm,scale=1.300]
  \draw[given] (0,0.75) -- (1,0.75);
  \dt{0,0.75} \dt{1,0.75}
  \node[vlab,above] at (0,0.75) {$A$};
  \node[vlab,above] at (1,0.75) {$B$};
  \draw[given] (0,0) -- (1,0);
  \foreach \x/\lb in {0/C,0.66/D,1/E}{%
    \dt{\x,0}
    \node[vlab,above] at (\x,0) {$\lb$};}
\end{tikzpicture}
\figcap{5--9}
\end{bkfigure}}

% --------------------------------------------------------------- Fig. 5-10
% Change of scale: E is the midpoint of the unit CD, so AB = 3 with unit CD
% and AB = 6 with unit ED.  Horizontal layout re-measured off the book plate
% (all in CD units from C): A at 1.55, B at 4.55, the half-unit ED alongside
% at 5.15..5.65.  The earlier 2.24/5.24/5.92/6.42 spread the plate too wide.
\newcommand{\FIGVTEN}{%
\begin{bkfigure}[0.96\linewidth]
\begin{tikzpicture}[x=1.25cm,y=1cm,scale=1.300]
  % unit CD with midpoint E, braced and labelled 1
  \draw[given] (0,0) -- (1,0);
  \dt{0,0} \dt{1,0}
  \draw[tick] (0.5,-0.09) -- (0.5,0.09);
  \node[vlab,below] at (0,-0.02) {$C$};
  \node[vlab,below] at (0.5,-0.02) {$E$};
  \node[vlab,below] at (1,-0.02) {$D$};
  \draw[fig,line width=0.5pt,decorate,decoration={brace,amplitude=3pt}]
        (0,0.16) -- (1,0.16);
  \node[slab,above,yshift=2.6pt] at (0.5,0.19) {1};
  % AB, three units long; the "1" sits under its first third
  \draw[given] (1.55,0.55) -- (4.55,0.55);
  \dt{1.55,0.55} \dt{4.55,0.55}
  \node[vlab,above] at (1.55,0.55) {$A$};
  \node[vlab,above] at (4.55,0.55) {$B$};
  \draw[tick] (2.55,0.46) -- (2.55,0.64);
  \draw[tick] (3.55,0.46) -- (3.55,0.64);
  \node[slab,below,yshift=-2.2pt] at (2.05,0.53) {1};
  % the half-unit ED alongside
  \draw[given] (5.15,0.30) -- (5.65,0.30);
  \dt{5.15,0.30} \dt{5.65,0.30}
  \node[vlab,above] at (5.15,0.30) {$E$};
  \node[vlab,above] at (5.65,0.30) {$D$};
\end{tikzpicture}
\figcap{5--10}
\end{bkfigure}}
